\subsection{Mutability and Identity: Immutable Objects, Mutable Objects, Shared References, and \texttt{is} vs. \texttt{==}}

Mutability asks whether an object's state can change in place while its identity---the address of its header on the heap---remains constant. Mutable cargo (lists, dicts, most user instances) supports in-place change; immutable cargo (ints, strings, tuples of immutables) requires rebinding to new objects to reflect new values.

Identity and equality answer different engineering questions:

\begin{itemize}
    \item \textbf{\texttt{is}}: compare pointer addresses of the wrapped objects---blindingly fast native pointer equality in CPython.
    \item \textbf{\texttt{==}}: invoke rich comparison (\texttt{\_\_eq\_\_}) that may traverse nested fields to decide value agreement.
\end{itemize}

Two distinct list objects with the same elements satisfy \texttt{==} but not \texttt{is}. Aliased names satisfy both. Immutable interning caches may make some equal ints share identity, but logic must never rely on that implementation detail.

Shared references to mutable cargo make aliasing bugs possible: mutation through one label visible through another until rebinding splits the arrows. Copy operations (\texttt{copy}, \texttt{list.copy}) create new outer containers; shallow copies still share inner references.

For C programmers: treat \texttt{is} like comparing addresses of structs; treat \texttt{==} like calling a typed comparator function that inspects contents.
